\section{Capítulo~\ref{sec:code}. Código Morse}

As estimativas-limite do capítulo são teoria da informação e contas; foram medidos onde surge o ruído no canal, quão rápido o cérebro trabalha e quanto custa um disparo, e a questão de que código o córtex usa de fato continua em aberto.

{\footnotesize
\setlength{\tabcolsep}{3pt}\renewcommand{\arraystretch}{1.15}
\begin{longtable}{>{\raggedright}p{0.5cm}>{\raggedright}p{4.0cm}>{\raggedright}p{5.6cm}>{\raggedright}p{2.3cm}>{\raggedright\arraybackslash}p{1.7cm}}
\toprule
N.º & Proposição & Experimento e o que foi medido & Fonte & Status \\
\midrule\endfirsthead
\toprule
N.º & Proposição & Experimento e o que foi medido & Fonte & Status \\
\midrule\endhead
1 & As frequências dos neurônios \nt{GLUT} do córtex vão de frações a dezenas de disparos por segundo, e a maior parte do tempo os neurônios trabalham perto do limite inferior & Registro com pipeta encostada à célula (a escolha do neurônio não depende de sua atividade) no córtex auditivo do rato acordado: a cada momento, menos de $5\,\%$ dos neurônios têm frequência acima de $20$ disparos por segundo; parte dos neurônios tem frequência de fundo abaixo de $0.01$ por segundo; a distribuição das frequências é lognormal & \cite{Hromadka2008} & medido \\
2 & A capacidade do canal de disparos é $R_{\max}\approx r\log_2\frac{e}{r\Delta t}$~\eqref{eq:spikeentropy}, e quase toda ela está no tempo dos disparos (proposição~\ref{prop:capacity}) & Entropia de uma cadeia binária de células de comprimento $\Delta t$. Números do capítulo: $8$ bits por disparo com $r=10$ por segundo e $\Delta t=1\ms$, cerca de $5$ bits com $\Delta t=10\ms$, menos de $2$ bits por segundo para um contador de disparos & \cite{MacKay1952} & teoria \\
3 & Os neurônios sensoriais transmitem de um a alguns bits por disparo, nos melhores casos até metade do limite & Reconstrução do estímulo a partir dos disparos de neurônios sensoriais para os quais o estímulo é conhecido; resumo das medições no livro & \cite{Rieke1997} & medido \\
4 & O gerador de disparos é preciso; o tempo exato é definido pelas variações rápidas da entrada & Neurônios do córtex de rato em fatias: com uma corrente repetida com flutuações rápidas, os disparos se repetem com precisão melhor que $1\ms$; com corrente constante, os instantes se espalham & \cite{Mainen1995} & medido \\
5 & A sinapse não é confiável: mais da metade dos disparos não chega ao destinatário por causa do acaso da liberação & Estimulação mínima das entradas de neurônios piramidais do hipocampo (fatias): mais da metade dos disparos não provoca resposta. Foram excluídas flutuações do limiar de estimulação, falhas de condução nas ramificações do axônio e a baixa temperatura do experimento; resta a liberação probabilística & \cite{AllenStevens1994} & medido \\
6 & O fluxo de disparos no córtex vivo parece aleatório: os intervalos se espalham como num fluxo de Poisson, e a variância do número de disparos fica próxima da média; parte do ``ruído'' é mensagem sobre os movimentos do animal & Registros nas áreas visuais V1 e MT do macaco acordado: coeficiente de variação dos intervalos de cerca de $1$, razão entre variância e média do número de disparos como num processo aleatório; um modelo que soma pequenas entradas independentes dá uma dispersão muito menor. No córtex visual do camundongo, uma parte considerável da variação é prevista pelo vídeo dos movimentos do próprio animal & \cite{Softky1993,Stringer2019} & medido \\
7 & O código de taxa é lento: erro $1/\sqrt{rT}$~\eqref{eq:rateerror}, enquanto o cérebro reconhece uma imagem em $\sim150\ms$ & A fórmula é estatística de Poisson, dedução do capítulo. Experimento: pessoas decidiam se havia um animal numa fotografia desconhecida mostrada por $20\ms$; os potenciais cerebrais para ``sim'' e ``não'' divergem após cerca de $150\ms$. A divisão desse tempo em uma dezena de estágios de $10$--$20\ms$ é estimativa do capítulo, não medição & \cite{Thorpe1996} & medido \\
8 & A média sobre um grupo é limitada pela correlação: o erro cai no máximo $1/\sqrt c$ vezes~\eqref{eq:correlated} (proposição~\ref{prop:pooling}) & Pares de neurônios da área MT do macaco registrados ao mesmo tempo: coeficiente de correlação dos números de disparos de $0.12$ em média, e a análise teórica mostra que mesmo essa correlação limita o ganho do grupo. Teoria: o limite é imposto só pela parte da correlação que se parece com o sinal. Modelo da posição oposta: a frequência de um grupo de $50$--$100$ neurônios é lida num único intervalo entre disparos, $10$--$50\ms$, e o tempo de disparos individuais não é usado no córtex & \cite{Zohary1994,MorenoBote2014,ShadlenNewsome1998} & medido \\
9 & O número pode ser escrito no tempo do primeiro disparo~\eqref{eq:latency}, na ordem dos disparos de um grupo ou na fase de um ritmo local & Retina da salamandra: a estrutura espacial de uma imagem mostrada brevemente está escrita nas latências relativas dos primeiros disparos dos neurônios de saída, quase independentemente do contraste. Células de lugar do hipocampo do rato: ao passar pelo ``seu'' lugar, os disparos se deslocam cada vez mais cedo no ciclo do ritmo teta ($7$--$12$\,Hz), um deslocamento de $100^\circ$ a $355^\circ$. Quanto o córtex usa o tempo dos disparos é questão em aberto (linha~8) & \cite{Gollisch2008,OKeefe1993} & medido \\
10 & Qual código é lido é decidido pela constante de tempo do destinatário~\eqref{eq:reader} (proposição~\ref{prop:reader}); no córtex ativo, os neurônios estão mais perto de detectores de coincidência & A fórmula~\eqref{eq:reader} é dedução do capítulo. Revisão de registros in vivo: no córtex ativo, a condutância da membrana é várias vezes maior que em repouso, e a constante de tempo, correspondentemente mais curta. Que o córtex troca de código exatamente assim não foi mostrado diretamente & \cite{Destexhe2003} & hipótese \\
11 & A sinapse com depressão transmite a mudança de frequência, essencialmente $\Delta r/r$~\eqref{eq:depression} (fig.~\ref{fig:depression}) & Registros pareados de neurônios piramidais do córtex: a velocidade da depressão é dada pela probabilidade de liberação; com depressão lenta, a resposta reflete a frequência; com rápida, as coincidências. Modelo baseado na depressão medida no córtex de rato: mudanças relativas iguais de frequência em entradas rápidas e lentas dão a mesma resposta, isto é, a sinapse transmite $\Delta r/r$. Os números $1.7\to2.2$, $6.7$ e $90\ms$ são cálculo do capítulo & \cite{Tsodyks1997,Abbott1997depression} & medido \\
12 & A rajada é o ``traço'': sua probabilidade de se perder, $(1-p)^k$~\eqref{eq:burst}, é pequena, e a facilitação a reduz ainda mais & Revisão: em sinapses pouco confiáveis, disparos isolados se perdem com frequência, e as rajadas são transmitidas de forma confiável graças à facilitação; a rajada provoca mudanças duradouras nas sinapses. Que o cérebro lê a rajada como um sinal à parte é hipótese & \cite{Lisman1997,AllenStevens1994} & hipótese \\
13 & Os neurônios \nt{GABA} rápidos definem o ritmo e a ``pontuação''; outra classe inibe os inibidores e abre a porta (proposição~\ref{prop:punctuation}) & Revisão das propriedades das classes de neurônios \nt{GABA} do córtex. Optogenética: a excitação rítmica de neurônios \nt{GABA} rápidos provoca o ritmo gama, e a resposta a um estímulo depende da fase do ciclo. No córtex visual do camundongo, os neurônios PV inibem uns aos outros, os SST inibem todas as outras classes, e os VIP inibem sobretudo os SST. A excitação dos neurônios VIP no camundongo acordado remove a inibição dos neurônios \nt{GLUT}; eles são ligados por um sinal de reforço. A divisão de papéis como regra geral é hipótese & \cite{Tremblay2016,Cardin2009,Pfeffer2013,Pi2013} & hipótese \\
14 & A energia limita a frequência média a algo da ordem de um disparo por segundo~\eqref{eq:energy} & Orçamento energético da substância cinzenta de roedores a partir de dados anatômicos e fisiológicos: disparos e correntes sinápticas respondem por $47\,\%$ e $34\,\%$ da energia de sinalização; no máximo $\sim15\,\%$ dos neurônios podem estar ativos ao mesmo tempo. Para o córtex humano: no máximo $\sim1\,\%$ dos neurônios podem estar fortemente ativos ao mesmo tempo. O número de $\sim10^9$ ATP por disparo e a potência de $\sim10$\,W para sinalização são estimativa do capítulo com base nesses cálculos & \cite{Attwell2001,Lennie2003} & teoria \\
15 & O código é esparso e ajustado para o máximo de bits por joule; o córtex troca precisão por energia (proposição~\ref{prop:economy}) & Hipótese do código econômico. Base: em camundongos com restrição alimentar, no córtex visual caem a condutância dos receptores AMPA e o gasto sináptico de ATP ($-29\,\%$), a frequência de disparos se mantém, e a sintonia à orientação se alarga $32\,\%$ & \cite{Barlow1961,Padamsey2022} & hipótese \\
\bottomrule
\end{longtable}}
